\section{Counterfactual Data Construction and Quality Control}
\label{app:counterfactual-data}

This appendix specifies how the source-free supervision was constructed and
frozen before training.  The procedure operates on registered event fields,
generates a new video from a source-free prompt, and screens that video without
access to a treatment label or any trained-model output.  It therefore defines
the training target; it does not measure the effect of \methodshort.  Digest
strings below are 12-hex SHA-256 prefixes; the cited registries store the full
digests.

\subsection{Structured Factual and Target Prompts}
\label{app:structured-prompts}

\paragraph{Candidate grid.}
For each mechanism, we form the Cartesian product of eight original training
sources, 12 training receivers, and two prompt styles (\emph{direct} and
\emph{natural}), giving \(8\times12\times2=192\) candidate bindings.  A
binding records the mechanism, source identity and phrase, receiver identity
and clean state, trigger, footprint, counterfactual state, prompt style, and
generation seed.  Factual and target prompts are rendered from these fields;
neither is obtained by deleting a noun from a free-form sentence.

For the six newly constructed mechanisms, the factual surface form is

\begin{quote}\footnotesize\ttfamily\raggedright
A simple realistic locked-camera video: <receiver> is visible;
<clean state>. <style-specific event clause>. Afterward, <footprint>;
<source end state>.
\end{quote}

The event clause contains the assigned source, the registered interaction,
and the receiver.  The direct and natural variants express the same field
binding with different syntax (for example, ``falls from above and visibly
strikes'' versus ``falls naturally onto and hits'').  The corresponding target
surface form is

\begin{quote}\footnotesize\ttfamily\raggedright
A simple realistic locked-camera close-up of <receiver>.
<counterfactual state>. <style-specific natural receiver motion>.
The viewpoint, receiver identity, lighting, and background geometry remain
stable.
\end{quote}

Water uses the previously frozen version of the same construction.  Its
factual template describes a source entering a named water-filled receiver,
the resulting splash and ripples, a moving reflected highlight, and stable
viewpoint and geometry.  Its target template retains the named receiver,
smooth water and gradual reflected-light motion, while omitting the entering
source, water contact, splash, and ripples.  The retained Water prompts are
validated byte-for-byte against their historical manifest rather than
regenerated from a revised template.

\paragraph{Registered examples.}
The following examples are exact prompt pairs from the candidate manifest
(line wrapping is typographic only).

\begin{quote}\small
\textbf{Direct, Water impact---factual.}
\emph{A simple realistic close-up video in one continuous shot. One large
clear water droplet falls naturally from above, enters the center of the water
in a calm shallow pond, and makes contact. The contact produces a visible
brief splash followed by circular ripples spreading across the water. A soft
reflected highlight moves slowly from left to right across the water and
receiver throughout the shot. The viewpoint, receiver, and background
geometry remain stable.}

\textbf{Direct, Water impact---target.}
\emph{A simple realistic close-up video in one continuous shot showing a calm
shallow pond. The water surface stays smooth, level, and undisturbed. A soft
reflected highlight moves slowly from left to right across the water and
receiver. The viewpoint, receiver, and background geometry remain stable.}

\textbf{Natural, Brittle fracture---factual.}
\emph{A simple realistic locked-camera video: one clear glass bowl on a dark
tabletop is visible; the brittle receiver is intact, uncracked, and still. As
the scene unfolds, one polished steel ball falls naturally onto one clear
glass bowl on a dark tabletop and hits its center. Afterward, the receiver
breaks and separate fragments remain around the impact point; the source rests
visibly among or beside the receiver.}

\textbf{Natural, Brittle fracture---target.}
\emph{A simple realistic locked-camera close-up of one clear glass bowl on a
dark tabletop. The brittle receiver remains intact and uncracked. Throughout
the shot, a reflected highlight moves gently across the intact surface. The
viewpoint, receiver identity, lighting, and background geometry remain
stable.}
\end{quote}

These pairs illustrate the intended textual intervention: the target retains
the receiver and low-amplitude scene dynamics but contains neither the source
nor language describing the registered trigger or footprint.  Mechanism-
specific forbidden-lexeme checks additionally reject target prompts that
contain the original source phrase or registered trigger/footprint terms.

\subsection{Dynamic Distributional Counterfactual Targets}
\label{app:distributional-targets}

Each target video is sampled directly from its source-free target prompt with
the frozen Wan2.1-T2V-1.3B generator.  Generation uses 25 inference steps,
guidance scale 5, bfloat16 inference, and an explicitly registered seed, and
produces 49 frames at 8 fps and \(832\times480\) resolution.  The
mechanism-specific exclusion phrase is supplied as negative conditioning; it
lists source-object cues, trigger, and footprint content forbidden for that
target.  Water contributes 192 previously generated, manifest-bound
candidates.  The other six mechanisms contribute 1,152 newly generated
candidates, for 1,344 candidates in total.  The new-generation aggregate
records all 1,152 expected videos as decoded and valid.

Crucially, a target is not produced by editing or inverting a factual video.
The factual prompt and target prompt share a registered receiver and context,
but the target is an independent draw from the generator.  Natural motion of
the receiver, illumination, or other non-event scene content is allowed.  The
construction consequently supplies a screened \emph{distributional
approximation} to the source-free world: it does not claim that every residual
cue is absent, or that object trajectories, textures, or pixels are aligned
with a particular factual realization. The candidate manifest is committed by
SHA-256
\texttt{80522812b906};
the validated new-target generation manifest is committed by
\texttt{b86fe17ea2ca}.

\subsection{Treatment-Blind Target Screening}
\label{app:target-screening}

Screening asks whether a candidate is usable as supervision, not whether a
method succeeds. For the six newly generated mechanisms, blind agents see the
complete 49-frame target through five
overlapping panels covering frames 0--12, 9--21, 18--30, 27--39, and 36--48.
Candidate identities are replaced by anonymous review identifiers and the
within-mechanism order is hash-randomized.  The review package explicitly
contains no treatment output or formal-evaluation video.

Each candidate from the six new mechanisms receives six atomic scores in
\(\{0,1,2\}\). The fields are receiver recognizability, source absence,
trigger absence, footprint absence, visual quality, and natural motion. Zero
denotes failed satisfaction---for example, an absent receiver, clearly
present forbidden content, or unusable quality; one denotes partial or
uncertain satisfaction, and two denotes clear satisfaction. A zero in any of
the first five fields is a hard rejection. Natural motion is a soft ranking
signal and never rejects one of these 1,152 new candidates by itself. The
survivors are ranked, in order, by the minimum and sum of the three core
absence fields, receiver recognizability, quality, natural motion, and the sum
of all six scores.  Exact ties are broken by ascending
\(
\operatorname{SHA256}(\texttt{target-select-v1}\mathbin{|}\texttt{candidate\_id})
\).

The six newly generated mechanisms received one full anonymous VLM review.
Four independently assigned blind agent-audit files then covered 239 targeted
candidates without access to the first pass's scores. Of 417 audited atomic
disagreements, 311 were resolved by two further blind adjudication passes that
did not receive either parent score.
The remaining 106 disagreements were retained as unresolved domains only
after exhaustive completion checks showed that neither endpoint could change
membership in the selected set. Water instead uses its earlier frozen,
named technical and semantic screen, which records 178 accepts and 14 rejects.
That legacy rule could reject a nearly static target on motion alone and must
not be conflated with the six-field rule above. Thus, all seven selections are
fixed without reference to either training arm.

The final treatment-blind selection manifest replaces the unused static
pre-media selection contract in the build registry. The initial contract
treated all six fields as eligibility conditions and proposed a hash-only
choice among candidates with fully absent core content. After anonymous target
review and targeted audit, but still before either training arm, the final
manifest froze the first-five-field hard rejection and score-based ranking
reported here, with natural motion used only in ranking. We therefore call
this the frozen treatment-blind selection rule, not a pre-registered method
outcome rule.

\begin{table}[t]
  \centering
  \caption{Treatment-blind target selection. \emph{Admissible} is the number
  remaining after the frozen hard-rejection rule; candidates above the quota
  are held back by the final frozen ranking rather than relabelled as failures.}
  \label{tab:target-screen-counts}
  \small
  \begin{tabular}{lrrr}
    \toprule
    Mechanism & Candidates & Admissible & Selected \\
    \midrule
    Water impact        & 192 & 178 & 178 \\
    Rigid collision     & 192 & 186 & 178 \\
    Brittle fracture    & 192 & 178 & 178 \\
    Powder impact       & 192 & 191 & 178 \\
    Elastic deformation & 192 & 184 & 178 \\
    Material release    & 192 & 179 & 178 \\
    Surface trace       & 192 & 191 & 178 \\
    \midrule
    Total               & 1,344 & 1,287 & 1,246 \\
    \bottomrule
  \end{tabular}
\end{table}

Observed hard-rejection reasons were confined to the registered fields:
missing or unrecognizable receivers, residual source cues, residual trigger or
footprint evidence, and unusable visual quality.  For example, the historical
Water rejects include an unrecognizable receiver, an entering hand or object,
surface patches resembling an unwanted footprint, and a nearly static target.
No candidate was rejected because a trained method later performed poorly.
If fewer than 178 candidates survive for a mechanism, the data version is
invalidated rather than replenished or screened after training.  The final
selected set contains exactly 178 targets per mechanism and is committed by
\texttt{d5d97dfb9868}.

\subsection{Source-Bank Separation and Assignment}
\label{app:source-bank}

Each mechanism has a 64-item training source bank comprising the eight
original training identities and 56 augmentation identities.  A separate set
of 48 identities is reserved for held-out-source evaluation.  The registry
requires the training bank and held-out set to be disjoint by source ID,
normalized phrase, and semantic family; held-out identities never enter a
training assignment.  For five mechanisms with compatible compact rigid
sources, the same public augmentation and holdout vocabularies are imported,
but receiver ontologies, prompts, compatibility tags, and adapters remain
mechanism-specific.  Material release and surface trace use their own
mechanism-compatible pointed and imprinting source banks.

For \methodshort, the 64-item bank is ordered by a deterministic SHA-256 rank
derived from the protocol, mechanism, and frozen ontology.  The resulting
permutation is repeated over the scheduled 100 erasure updates and then over
the 78 inactive selected rows.  Within each partition, deterministic swaps
repair any collision with that row's original source.  Consequently, the
assigned source differs from the original in both registered ID and normalized
phrase for every one of the 178 rows.  This is the sense in which the original
identity is excluded from an assignment; the other seven original-training
identities remain legitimate members of the 64-item bank.

All 64 identities occur in the active 100-row schedule: 36 occur twice and 28
once. Across the complete 178-row mapping, 50 identities occur three times
and 14 twice. These counts hold for every mechanism. The complete 178-row
assignment is stored in a registry-bound CSV. The per-mechanism registry
commits to a second hash of the deterministically derived assignment salt and
to the active-100 and full-178 projections. The ontology
registry is committed by
\texttt{37ef4920665e},
and the top-level training-input registry binds all seven mapping registries.

\subsection{Preservation Bank and Cache Integrity}
\label{app:preservation-data}

\paragraph{Shared preservation bindings.}
All 18 Wan method/control adapters use the same 36 generic preservation
bindings.  They consist of
12 scene contents---books, vase, lamp, clock, fan, toy train, pendulum,
conveyor, walking person, clouds, traffic light, and turntable---each rendered
with three registered prompt variants.  The variants respectively request
stable identities and context, exclude abrupt target-like physical events,
or exclude entries from above and abrupt events.  Each prompt is paired with
a separately generated 49-frame target video.  The frozen preservation CSV is
committed by
\texttt{d5cbe96ac838}
and is shared byte-for-byte by all seven mechanism registries.

\paragraph{What is cached.}
For each mechanism, the base cache has 214 entries: 178 selected target-video
latents with their canonical factual-prompt embeddings, followed by 36
preservation latents and prompt embeddings.  The teacher cache contains the
178 source-free target-prompt embeddings, and each arm cache contains the 178
student-prompt embeddings for that arm.  Frozen-teacher predictions are
computed online at the sampled noise level with the adapter disabled; they are
not stored as target videos or precomputed output tensors.  All cached tensors
use \texttt{torch.bfloat16}.  Video latents have shape
\([1,16,13,60,104]\), and prompt embeddings have shape
\([1,226,4096]\).

\paragraph{Integrity contract.}
Before encoding, every selected and preservation video is decoded and checked
for one video stream, no audio, 49 frames, 8 fps, and \(832\times480\)
resolution, and its bytes must match the registered SHA-256.  Each cache entry
then records the SHA-256 of the serialized file and of each latent or prompt-
embedding tensor, including its shape and dtype in the tensor digest.  Each
cache manifest additionally binds the ordered file inventory, selected-target
CSV, preservation CSV, model inventory, runtime registry, and cache-input
registry. Before model initialization, training validates the frozen run
specification and cache manifests as well as every loaded tensor's identity,
shape, dtype, and finiteness. It repeats the frozen registry and manifest
binding checks after optimization and before checkpointing.

The model is the frozen Wan2.1-T2V-1.3B Diffusers checkpoint.  Its content
inventory contains 94 files and 28,935,657,115 bytes, with ordered content
digest
\texttt{717d6d2d9518};
the inventory file itself has digest
\texttt{51e7199b99ee}.
The frozen runtime registry has digest
\texttt{9043adf7f823}.
Finally, each of the seven Matched-Control receipts compares all 178 factual
prompt-embedding tensors against the corresponding base-cache tensors and
requires exact equality.  The run-spec registry was released only after all
cache manifests and equality receipts passed; it is committed by
\texttt{36c79a862918}.
